\documentclass[10pt,a4paper]{amsart}
\usepackage[margin=19mm]{geometry}
\usepackage{amsmath,amssymb,mathtools}
\usepackage[scr=rsfs,cal=euler]{mathalfa}
\usepackage{xfrac}
\usepackage[unicode,hidelinks,bookmarks=false]{hyperref}
\newcommand{\ie}{i.e.\ }
\newcommand{\cf}{cf.\ }
\newcommand{\resp}{respectively\ }
\newcommand{\Subsection}{\subsection}
\providecommand{\nobreakdash}{\nobreak\hbox{-}}
\setlength{\emergencystretch}{2em}
\multlinegap0pt
\usepackage{bm}
\usepackage{bbm}
\usepackage{dsfont}
\usepackage{xspace}
\usepackage{cancel}
\usepackage{mathtools}
\usepackage{amsthm}
\makeatletter
\@ifundefined{theorem}{\newtheorem{theorem}{Theorem}}{}
\makeatother
\title[Gauges and Accelerated Optimization over Smooth and/or Stron]{Gauges and Accelerated Optimization over Smooth and/or Strongly Convex Sets}
\author{Ning Liu, Benjamin Grimmer}
\date{Source: arXiv:2303.05037v4 (CC BY 4.0)}
\begin{document}
\maketitle

\noindent\emph{Source and licence.} Gauges and Accelerated Optimization over Smooth and/or Strongly Convex Sets. Version arXiv:2303.05037v4 (CC BY 4.0), distributed under the Creative Commons Attribution 4.0 International licence, as recorded in the arXiv licence metadata for this exact version and shown on the abstract page https://arxiv.org/abs/2303.05037v4. Full rights notice: licenses/arxiv-2303.05037-CC-BY-4.0.txt.\par\smallskip

\noindent\textit{Editorial scope.} Counted unit: the Theorem whose source label is \texttt{thm:level-proj-convergence} in the article (source0.tex lines 2250--2256, source label \texttt{thm:level-proj-convergence}), delivered together with the article's own complete original proof of it (source0.tex lines 2257--2282, one \texttt{proof} environment of 26 lines). No mathematics is written, repaired or re-derived: statement and proof are the article's own text line for line. Required context retained uncounted: eq:general-subgrad-bound (lines 2122--2124). Every definition, display and label of those lines is retained unchanged. Omitted from this package and not redistributed: 1--2121, 2125--2249, 2283--2385. Nothing inside a retained excerpt is omitted or altered: every retained body is delivered exactly as the article's lines are written, so no omission receipt of any kind accompanies this package. Citations: the retained excerpts carry no citation command at all, so no bibliography entry of the article is transcribed and no reference is redistributed. Reference adaptation: none was needed and none was made; every reference inside the delivered unit resolves inside it, so no unresolved reference or citation is produced. Printed slips kept literal: no printed word, symbol, hypothesis or number of the counted statement or of its proof was altered. Layout and class: the article's own class is \texttt{sn-jnl} with packages including graphicx, multirow, amsmath, amssymb, amsfonts, amsthm, mathrsfs, appendix, xcolor, textcomp, manyfoot, booktabs, algorithm, algorithmicx; the wrapper is the lane's pinned \texttt{amsart} editorial wrapper at 10pt on A4, which restates the theorem-like environments the retained text needs and adds the article's own macro definitions that the retained text needs (none), with extra wrapper packages (bm, bbm, dsfont, xspace, cancel, mathtools). The delivered numbering is the unit's own. Assets: no retained span contains a figure environment, a graphics-inclusion command, a PDF-page inclusion, a table or a TikZ picture; no image file is copied into this package, the package declares no asset dependency and no image file is redistributed with it. Licence: version arXiv:2303.05037v4 of this article is distributed under the Creative Commons Attribution 4.0 International licence, as recorded in the arXiv licence metadata for this exact version and shown on the abstract page https://arxiv.org/abs/2303.05037v4. A full rights notice is delivered at licenses/arxiv-2303.05037-CC-BY-4.0.txt. Characters: every retained excerpt is pure ASCII, so the delivered engine, XeLaTeX with its default Latin Modern text font, reports no missing character.
\begin{equation}
    \|h_i(y)g_i\|^2 \leq M^2h_i^2(y) = M^2p_*^2 + 2M^2\left(\frac{1}{2}h^2_i(y) - \frac{1}{2}p_*^2\right). \label{eq:general-subgrad-bound}
\end{equation}

\begin{theorem}\label{thm:level-proj-convergence}
    For any convex, nonnegative, $M$-Lipschitz functions $h_i$ and $\bar h \geq p_*$, the level-set projection method $y_{k+1}=\mathtt{level\mbox{-}proj}(y_k,\bar h)$  has
    $$ \min_{k=0\dots T}\left\{h(y_k) - \bar h\right\} \leq \frac{MD}{\sqrt{T+1}} + \frac{2M^2D^2}{\bar h(T+1)} $$
    provided some $\bar y$ with $f(\bar y)\leq \bar h$ has $\|y_0-\bar y\|\leq D$.
    If additionally each $\frac{1}{2}h_i^2$ is $\mu$-strongly convex and $T+1 \geq \frac{32M^2}{\mu}\log(\mu D^2/\bar h^2)$, then
    $$ \min_{k\leq T}\{h(y_k) - \bar h\} \leq \frac{\sqrt{32}M^2\bar h}{\mu (T+1)} + \frac{64M^4\bar h}{\mu^2(T+1)^2} \ . $$
\end{theorem}

\begin{proof}
Let $\mathcal{\bar Y} = \{y \mid \frac{1}{2}h^2(y) \leq \frac{1}{2}\bar h^2\}\neq\emptyset$ denote the target level set, and $\mathcal{Y}_k = \{y \mid m_k(y)\leq \frac{1}{2}\bar h^2\}\supseteq \mathcal{\bar Y}$ denote our model's target level set (the containment follows from $m_k$ lower bounding $\frac{1}{2}h^2$). Each orthogonal projection step ensures $y_{k}-y_{k+1}$ is normal to $\mathcal{Y}_k$ at $y_{k+1}$. Consequently every $\bar y\in \mathcal{\bar Y}\subseteq\mathcal{Y}_k$ must have $(y_k - y_{k+1})^T(y_{k+1}-\bar y)\geq 0$. Then
\begin{align}
    \|y_{k+1}-\bar y\|^2 &= \|y_k - \bar y\|^2  -  2(y_{k} - y_{k+1})^T(y_{k+1}-\bar y) - \|y_{k+1}-y_k\|^2 \nonumber \\
    &\leq \|y_k - \bar y\|^2 - \|y_{k+1}-y_k\|^2 \nonumber \\
    &\leq \|y_k - \bar y\|^2 - \frac{\left(\frac{1}{2}h^2(y_k) - \frac{1}{2}\bar h^2\right)^2}{M^2_k} \nonumber  \\
    &= \|y_k - \bar y\|^2 - \frac{\left(\frac{1}{2}h^2(y_k) - \frac{1}{2}\bar h^2\right)^2}{M^2\bar h^2 + 2M^2\left(\frac{1}{2}h^2(y_k) - \frac{1}{2}\bar h^2\right)} \ . \label{eq:level-proj-induction}
\end{align}
where the last inequality follows as $m_k$ is $M_k$-Lipschitz (by~\eqref{eq:general-subgrad-bound}) with $m_k(y_k)-m_k(y_{k+1}) = \frac{1}{2}h^2(y_k) - \frac{1}{2}\bar h^2$, and so $\|y_{k+1}-y_k\| \geq (\frac{1}{2}h^2(y_k) - \frac{1}{2}\bar h^2)/M_k$. Inductively applying~\eqref{eq:level-proj-induction} yields
\begin{align*}
    & \min_{k=0\dots T}\left\{\frac{\left(\frac{1}{2}h^2(y_k) - \frac{1}{2}\bar h^2\right)^2}{M^2\bar h^2 + 2M^2\left(\frac{1}{2}h^2(y_k) - \frac{1}{2}\bar h^2\right)}\right\} \leq \frac{\|y_0-\bar y\|^2}{T+1} \\
    \implies & \min_{k=0\dots T}\left\{\frac{1}{2}h^2(y_k) - \frac{1}{2}\bar h^2\right\} \leq \frac{M\bar h\|y_0-\bar y\|}{\sqrt{T+1}} + \frac{2M^2\|y_0-\bar y\|^2}{T+1} \ . 
\end{align*}

Supposing $\mu$-strong convexity holds, observe that for $\bar y\in \mathcal{\bar Y}$ closest to $y_k$, $\frac{1}{2}h^2(y_k) - \frac{1}{2}\bar h^2 \geq \frac{\mu}{2}\|y_k-\bar y\|^2$. Then letting $D_k = \mathrm{dist}(y_k, \mathcal{\bar Y})^2$,~\eqref{eq:level-proj-induction} can then be relaxed to the recurrence relation
$$ D_{k+1} \leq D_k - \frac{\mu^2 D_k^2}{4(M^2\bar h^2 + \mu M^2D_k)} \ .$$
Note that this recurrence implies $D_k$ is monotonically decreasing. Let $K$ denote the first index with $D_k \leq \bar h^2/\mu$. All $k<K$ must have a geometric decrease of $D_{k+1}\leq (1-\mu/8M^2 )D_k$. Thus $K\leq (8M^2/\mu) \log(\mu D_0/\bar h^2)$. For all $k\geq K$, the recurrence ensures $D_{k+1}\leq D_k - \mu^2 D_k^2/8M^2\bar h^2$, which implies $D_k \leq 8M^2\bar h^2/\mu^2(k-K)$. \footnote{The recurrence $D_{k+1} \leq D_k - \alpha D_k^2$, is equivalent to $\frac{1}{D_k} \leq \frac{1}{D_{k+1}}-\alpha \frac{D_k}{D_{k+1}}$, then by $\frac{D_k}{D_{k+1}}\geq 1$ and $1/D_K>0$, we obtain $D_k$.}
By assumption, we have $T/2 > 2K$ and so $T/2-K+1 \geq T/4 +1 $. Hence recurrence ensures
$$ D_{T/2} \leq \frac{32 M^2\bar h^2}{\mu^2(T+1)} \ . $$
Then inductively applying~\eqref{eq:level-proj-induction} from $k=T/2$ to $T$ yields our claimed rate as
\begin{align*}
    & \min_{k=T/2\dots T}\left\{\frac{\left(\frac{1}{2}h^2(y_k) - \frac{1}{2}\bar h^2\right)^2}{M^2\bar h^2 + 2M^2\left(\frac{1}{2}h^2(y_k) - \frac{1}{2}\bar h^2\right)}\right\} \leq \frac{D_{T/2}}{T/2} \\
    \implies & \min_{k=0\dots T}\left\{\frac{1}{2}h^2(y_k) - \frac{1}{2}\bar h^2\right\} \leq \frac{M\bar h\sqrt{D_{T/2}}}{\sqrt{T+1}} + \frac{2M^2D_{T/2}}{T+1} \\
    \implies & \min_{k=0\dots T}\left\{\frac{1}{2}h^2(y_k) - \frac{1}{2}\bar h^2\right\} \leq \frac{\sqrt{32}M^2\bar h^2}{\mu (T+1)} + \frac{64M^4\bar h^2}{\mu^2(T+1)^2} \ . 
\end{align*}
\end{proof}

\end{document}
